\documentclass[11pt,letterpaper]{article}

\usepackage[margin=1in]{geometry}
\usepackage{amsmath,amssymb,amsfonts}
\usepackage{booktabs}
\usepackage{longtable}
\usepackage{array}
\usepackage{tabularx}
\usepackage{enumitem}
\usepackage{hyperref}
\usepackage{xcolor}
\usepackage{setspace}
\usepackage{fancyhdr}
\usepackage{caption}

\hypersetup{
  colorlinks=true,
  linkcolor=blue!50!black,
  urlcolor=blue!50!black,
  citecolor=blue!50!black
}

\setstretch{1.12}
\setlist[itemize]{leftmargin=1.4em,itemsep=0.25em}
\setlist[enumerate]{leftmargin=1.6em,itemsep=0.25em}
\setlength{\parskip}{0.35em}
\setlength{\parindent}{0pt}

\pagestyle{fancy}
\fancyhf{}
\lhead{Solar UAV optimizer}
\rhead{How the model actually works}
\cfoot{\thepage}

\newcommand{\flag}{\textcolor{red!60!black}{\textbf{FLAGGED}}}
\newcommand{\code}[1]{\texttt{#1}}

\newenvironment{intuition}{\begin{quote}\small}{\end{quote}}
\newenvironment{examplebox}{\begin{quote}\small}{\end{quote}}
\newenvironment{warningbox}{\begin{quote}\small}{\end{quote}}

\title{\textbf{How the Solar UAV Optimizer Works}\\[0.4em]
\large Assumptions, constraints, physics, and what the numbers are trying to tell us}
\author{Design toolchain notes for new members\\
\texttt{solar-uav-design} repository}
\date{17 August 2026}

\begin{document}
\maketitle
\tableofcontents
\newpage

% =====================================================================
\section{Start here: the one-page mental model}
% =====================================================================

This repository is not a CAD package and it is not a gradient-based ``optimal control'' solver.
It is a \textbf{physics model plus a discrete search}. For each candidate airplane it answers one question:

\begin{quote}
\textit{If this airframe flies a clear-sky solstice day at our site, with our packs and our motors, does the battery stay above 20\% overnight and refill by the next afternoon?}
\end{quote}

If yes, the design \emph{closes}. If not, the model reports how far it missed, in watt-hours. Nothing on the locked grid currently closes.

The night, not the day, is the problem. On a June solstice the array makes on the order of 200\,W at noon. Cruise only needs about 75\,W of bus power. The airplane still fails because for about \textbf{14 hours} the array cannot cover cruise, and two packs only store 773\,Wh of usable energy.

\begin{intuition}
\textbf{Intuition.} Think of the battery as a water tank, the motors as a tap, and the solar array as a daytime hose. Closing is not ``does the hose fill the tank?'' Closing is ``can the tank last until the hose starts working again, and then refill before the next night?'' Extra noon solar that you cannot store does not help the night.
\end{intuition}

Three numbers dominate almost every trade:

\begin{enumerate}
  \item \textbf{Night bus power} --- electrical watts the two motors plus avionics draw in level night loiter.
  \item \textbf{Battery-true night length} --- hours when array power is less than that night load (not sunset-to-sunrise).
  \item \textbf{Usable pack energy} --- $2 \times 483\,\mathrm{Wh} \times 80\% = 773\,\mathrm{Wh}$.
\end{enumerate}

Aero, mass, props, and motors all matter only insofar as they move those three numbers.

Source of truth for locked constants: \code{solar\_uav/config.py}. Every value marked \flag{} is an assumption awaiting a measurement or a team decision. The physics rule is: when uncertain, take the \textbf{pessimistic} side. A design must not close because the model was kind.

% =====================================================================
\section{The airplane and the mission}
% =====================================================================

\subsection{What we are designing}

Twin-boom, OV-10-style solar UAV:

\begin{itemize}
  \item One three-piece wing (S4310): rectangular inboard bay between the booms, two outboard panels that may taper. Leading edge unswept.
  \item Independent tails: NACA~0010 H-stab spanning the booms \emph{in the wing plane}; two NACA~0010 V-stabs further aft (fin LE behind H-stab TE). Not a $\pi$-tail.
  \item One tractor propeller and one motor on each boom.
  \item SunPower C60 cells on the wing top and on the H-stab (ahead of the elevator).
  \item Two Upgrade Energy GOLD V1 6S packs in parallel, Genasun GVB-8-Li-25.0V boost MPPTs (one per string).
\end{itemize}

No payload. The mission is an 81-hour endurance record, designed to 89 hours (10\% margin). The optimizer itself simulates \textbf{one 24-hour design day} and requires the pack to refill, because an 89-hour flight is just that day repeated. If you end the afternoon at 95\% instead of 100\%, the next night starts lower, and the mission ratchets empty.

\subsection{Locked mission requirements}

\begin{center}
\begin{tabular}{@{}ll@{}}
\toprule
Quantity & Value \\
\midrule
Site latitude & $34.646849^\circ\mathrm{N}$ (El Mirage Race Course) \\
Site longitude & $-117.605977^\circ$ (El Mirage Race Course) \\
Site elevation & $2837\,\mathrm{ft} = 864.7\,\mathrm{m}$ MSL (USGS lake bed) \\
Cruise AGL & $400\,\mathrm{ft}$ (density at field + AGL; no ground-effect credit) \\
Timezone & America/Los\_Angeles \\
Design day & 21 June 2027 (solstice, longest day) \\
Mission window & 1 May -- 31 July 2027 \\
MTOW & $\le 12\,\mathrm{kg}$ \\
Span & $\le 6\,\mathrm{m}$ \\
Packs & at most three (integer; 2 or 3 searched) \\
Avionics & $5\,\mathrm{W}$ continuous, day and night (measured) \\
Climb & $\ge 1.5\,\mathrm{m/s}$ \flag{} \\
Min airspeed & $\max(1.2\,V_\mathrm{stall},\; 15\,\mathrm{kt})$ \\
SOC window & 20\%--100\% (80\% of nameplate is usable) \\
\bottomrule
\end{tabular}
\end{center}

Temperatures are a \textbf{hot design-day}: ambient $24^\circ\mathrm{C}$ dawn to $40^\circ\mathrm{C}$ afternoon \flag{}. A real El Mirage June dawn is often cooler and denser. Keeping $24^\circ\mathrm{C}$ dawn is conservative for night aero. Afternoon $40^\circ\mathrm{C}$ is not a desert extreme --- revisit if cells run hotter than the 75\,$^\circ\mathrm{C}$ noon calibration.

% =====================================================================
\section{The conservative-physics rule}
% =====================================================================

If a number is not measured and not in a datasheet range, the model implements the side that makes closure \emph{harder}.

\begin{itemize}
  \item \textbf{Prop maps.} If the (speed, RPM) point is on the manufacturer table, use that $T$ and $P$, density-scaled. Similarity $C_T(J)$ is a check, or a fallback only while the blade is still below windmill $J$. Never invent a slower RPM at fixed speed once a tabulated block already covers that speed.
  \item \textbf{ESC / motor.} Do not treat 95\% ESC efficiency as valid below $\sim$50\% duty. Low-RPM cruise is not a free no-load-current win.
  \item \textbf{Aero.} Do not take unmodeled credits (T-tail height reducing downwash, flight cooling, swirl recovery) unless they \emph{hurt} the metric you care about. Handbook downwash $2C_L/(\pi\,\mathrm{AR})$ is already coarse --- do not ``improve'' it in the direction that lightens tail download.
  \item \textbf{Mass / solar / battery.} Heavier structure, less array, stricter charge limits --- not the reverse.
\end{itemize}

\begin{warningbox}
\textbf{Why this exists.} An earlier interpolator $n^2$-scaled a $22\times12$E propeller down to $\sim$1190\,RPM at night cruise, inventing $J \approx 1.14$ past the map's zero-thrust $J \approx 0.67$. Night bus looked like $\sim$62\,W and the airplane ``closed.'' The honest operating point on the table is $\sim$2070\,RPM and $\sim$133--138\,W. Prefer an honest fail over an optimistic close.
\end{warningbox}

% =====================================================================
\section{What kind of optimizer this is}
% =====================================================================

File: \code{solar\_uav/optimize.py}.

It is a \textbf{mixed continuous / discrete search}. Aero geometry (span, chord, H-stab arm, V-stab arm, boom spacing, taper, washout, elevator fraction) is continuous and is driven by differential evolution (\code{search\_continuous()}). Integer and catalog quantities stay discrete: cell count from packing, pack count (2 or 3, cap 3), motor (A40-12L direct), ESC (fixed mass), and the inner-loop propeller pick.

A legacy exhaustive grid (\code{search()}) remains for neighborhood dumps. Do not use it as the production path --- it is the slow nested loop that used to take tens of minutes.

The inner loop (\code{\_best\_prop}) is not part of the outer product. For each airframe it:

\begin{enumerate}
  \item Tries every shortlisted propeller.
  \item Keeps only those that can make night thrust on 6S.
  \item Ranks them by night bus power (lowest first).
  \item Accepts the first that also meets the $1.5\,\mathrm{m/s}$ climb rule.
\end{enumerate}

So ``the optimizer evaluated 51 props'' does \emph{not} mean $51\times$ the airframe grid. Each airplane picks one fan.

Ranking of results: closed designs first, then highest \code{margin\_wh}. If none closed, the top row is the nearest miss.

Default motor is the A40-12L kv410 direct-drive. Geared A40-10L / 10S stay in the catalog for compare scripts only and are not searched.

% =====================================================================
\section{Optimization objective}
% =====================================================================

\textbf{Maximize 24-hour energy margin} (watt-hours above the 20\% SOC floor), subject to refill, climb, MTOW, span, solar layout, and trim/pitch.

After the time march,

\begin{align}
  E_\mathrm{reserve}
    &= (\mathrm{SOC}_\mathrm{min} - 0.20)\,E_\mathrm{pack,total}, \\
  E_\mathrm{refill}
    &= (\mathrm{SOC}_\mathrm{end} - 0.98)\,E_\mathrm{pack,total}, \\
  \mathrm{margin}
    &= \min(E_\mathrm{reserve},\, E_\mathrm{refill}).
\end{align}

If the pack could not supply the requested power at any step, margin is also capped at $-\,E_\mathrm{unmet}$.

\begin{intuition}
\textbf{Intuition.} $\mathrm{SOC}_\mathrm{min}$ is ``did we scrape the floor at dawn?'' $\mathrm{SOC}_\mathrm{end}$ is ``are we full again this afternoon?'' Taking the minimum of those two watt-hour gaps means a design that barely survives the night but cannot refill is not allowed to look like a winner. For an 89-hour flight, failing to refill is just a slow way to fail on night two.
\end{intuition}

A design is \textbf{closed} only if all of these hold:

\begin{itemize}
  \item $\mathrm{SOC}_\mathrm{min} \ge 0.20$
  \item unmet energy $< 0.5\,\mathrm{Wh}$
  \item climb $\ge 1.5\,\mathrm{m/s}$
  \item $\mathrm{SOC}_\mathrm{end} \ge 0.98$
\end{itemize}

Margin of $-158\,\mathrm{Wh}$ on the current nearest-miss means: at the worst point the pack is 158\,Wh below the 20\% floor (here $\mathrm{SOC}_\mathrm{min} = 0.036$). It is not ``158\,Wh of extra solar would close it,'' because extra noon solar mostly spills. You need that energy \emph{during the night}, which means lower night watts, a longer pack, or array power that turns on earlier in the morning. The miss is the 20\% reserve, not unmet energy (unmet is 0\,Wh); refill also fails ($\mathrm{SOC}_\mathrm{end} = 0.887$).

% =====================================================================
\section{Constraints}
% =====================================================================

\subsection{Hard constraints (candidate is discarded before the mission march)}

Implemented in \code{\_hard\_constraints} and the packing checks:

\begin{center}
\begin{tabular}{@{}lp{0.62\textwidth}@{}}
\toprule
Name & Rule \\
\midrule
Span & $b \le 6\,\mathrm{m}$ \\
Mass & $m \le 12\,\mathrm{kg}$ (re-checked after the prop is chosen) \\
Solar layout & cells fit the bays; each series string has $1$--$26$ cells; Voc-legal at cold cells \\
Stall & $1.2\,V_\mathrm{stall} \le 25\,\mathrm{m/s}$ (sanity cap) \\
Stability & static margin $\ge 8\%\,\mathrm{MAC}$, elevator trim, pitch acceleration $\ge 0.30\,\mathrm{rad/s^2}$, roll rate $\ge 35^\circ/\mathrm{s}$, differential-thrust yaw vs $5^\circ$ sideslip and yaw rate $\ge 5^\circ/\mathrm{s}$ (no rudder). $C_{n_\beta}$ / $|C_{n_r}|$ floors off; derivatives still computed. \\
Boom & carbon tube length $> 0.4\,\mathrm{m}$ \\
V-stab / H-stab & V-stab leading edge at or behind H-stab trailing edge (also enforced by flooring $\ell_v$) \\
Propulsion & some shortlisted prop can make night thrust on 6S \emph{and} climb \\
\bottomrule
\end{tabular}
\end{center}

\subsection{Energy constraints (scored, not discarded)}

Candidates that fail SOC or refill are still written to the CSV so we can see nearest misses. Climb failure also appears in the reason string.

\subsection{Electrical / packing constraints baked into geometry}

\begin{itemize}
  \item One series string per bay on the locked search (inboard, left outboard, right outboard, H-stab). Strings do not cross joints.
  \item String $V_\mathrm{oc}$ at $10^\circ\mathrm{C}$ cells must stay $\le 19.2\,\mathrm{V}$ (depleted 6S). That caps a string at \textbf{26 cells}.
  \item Genasun input $\le 8\,\mathrm{A}$; C60 $I_\mathrm{mp} \approx 5.9\,\mathrm{A}$, so one series string per MPPT is fine.
  \item Layout pitch $130\,\mathrm{mm}$ (cell is $125\,\mathrm{mm}$ plus gap). A $300\,\mathrm{mm}$ chord still fits two rows.
  \item After packing, every claimed cell is placed as a $125\,\mathrm{mm}$ square and rejected if it overlaps another cell or hangs off the planform (taper, aileron keep-out, elevator keep-out).
  \item H-stab cells only on the fixed portion; elevator chord is keep-out (split elevons).
  \item Tapered wing: no cells on the aileron chord. Rectangular wing: cells may cover ailerons.
\end{itemize}

% =====================================================================
\section{Decision variables and search grids}
% =====================================================================

Aero variables are continuous. Integer / catalog items stay discrete. Every continuous quantity has a \texttt{(low, high, start)} triple in \code{config.CONTINUOUS}. User-locked starts (chord $310\,\mathrm{mm}$, tail arm $0.90\,\mathrm{m}$, airspeed $12\,\mathrm{m/s}$) are not overwritten.

\begin{center}
\begin{tabular}{@{}lll@{}}
\toprule
Variable & Bounds / start & Notes \\
\midrule
\multicolumn{3}{l}{\emph{In the DE vector}} \\
Span $b$ & $4.50$--$6.00\,\mathrm{m}$, start $6.00$ & hard cap $6\,\mathrm{m}$ \\
Root chord $c_r$ & $0.300$--$0.52\,\mathrm{m}$, start $310\,\mathrm{mm}$ & two C60 rows; high = four pitches \\
H-stab arm $\ell_t$ & $0.80$--$2.00\,\mathrm{m}$, start $0.90$ & wing AC to H-stab AC \\
V-stab arm $\ell_v$ & $0.80$--$3.00\,\mathrm{m}$, start $1.20$ & floored so V-stab LE $\ge$ H-stab TE \\
Elevator fraction & $0.22$--$0.40$, start $0.34$ & split elevons; keep-out for H-stab cells \\
Taper $\lambda = c_\mathrm{tip}/c_r$ & $0.40$--$1.0$, start $0.60$ & $1.0$ is rectangular; $0.40$ still packs a tip cell \\
Taper start & $0$--$1$ of outboard, start $0.50$ & $0$ = taper from the boom \\
Boom spacing & $1.04$--$2.21\,\mathrm{m}$, start $1.30$ & $8$--$17$ cell pitches; H-stab span = boom spacing \\
Washout tip & $0^\circ$--$5^\circ$, start $0^\circ$ & $5^\circ$ is a lot at this AR \\
Washout start & $0$--$1$ of outboard, start $0$ & unused at $0^\circ$ washout \\
\midrule
\multicolumn{3}{l}{\emph{Continuous, nested (not DE)}} \\
Loiter airspeed & $8$--$18\,\mathrm{m/s}$, start $12$, $0.01\,\mathrm{m/s}$ step & still floored by $1.3\,V_s$ / $15\,\mathrm{kt}$ \\
\midrule
\multicolumn{3}{l}{\emph{Continuous, held at start}} \\
H-stab AR & $3.5$--$6.0$, start $4.5$ & not in the DE vector \\
V-stab AR & $1.5$--$3.5$, start $1.5$ & per fin; held at start \\
Aileron chord frac & $0.20$--$0.30$, start $0.25$ & locked for S\&C sizing \\
\midrule
\multicolumn{3}{l}{\emph{Discrete / catalog}} \\
Packs & $\{2,3\}$ & cap 3; one DE per pack count \\
Motor & A40-12L kv410 direct & not searched \\
ESC & fixed $0.060\,\mathrm{kg}$ lump & not searched \\
$N_\mathrm{cells}$ & from packing & integer, Voc-capped strings \\
Prop & inner loop, see \S\ref{sec:props} & discrete catalog \\
\bottomrule
\end{tabular}
\end{center}

The production search is \code{optimize.search\_continuous()} (differential evolution, start at the DE table above). A neighborhood grid dump is optional via the legacy \code{search()} path.

% =====================================================================
\section{How one candidate is scored}
% =====================================================================

Walk this loop in your head for every row of the CSV.

\begin{enumerate}
  \item Build geometry: planform, boom spacing, H-stab in the wing plane sized to $V_H = 0.50$ (chord also floored by one cell row), V-stabs to $V_V = 0.030$ at their own arm $\ell_v$ (no 1.30\,m area floor), carbon tube from wing $c/4$ to V-stab trailing edge, fuselage pods from wing TE to the prop disk.
  \item Pack cells: one Voc-legal series string per bay. Reject if cells overlap or hang off the planform.
  \item Sum mass. Drop if $>12\,\mathrm{kg}$.
  \item Check stall, static margin, elevator, pitch, roll rate, weathercock, yaw damping.
  \item For each propeller: night drag at min-power speed $\to$ RPM on the map $\to$ bus watts. Rank by night watts. Keep the cheapest that still climbs at 25.2\,V.
  \item March 24 hours: solar from pvlib, load = night bus or day bus, $P_\mathrm{net} = P_\mathrm{solar} - P_\mathrm{load}$, battery with 6\,A/pack charge cap.
  \item Record $\mathrm{SOC}_\mathrm{min}$, $\mathrm{SOC}_\mathrm{end}$, margin, flags (especially prop table vs.\ similarity).
\end{enumerate}

The rest of this document is the physics inside those steps.

% =====================================================================
\section{Atmosphere and site}
% =====================================================================

Files: \code{solar\_uav/aircraft.py} (density), \code{solar\_uav/environment.py} (sun and temperature).

Site pressure is ISA-style at field elevation plus $400\,\mathrm{ft}$ AGL ($h \approx 987\,\mathrm{m}$ MSL):

\begin{equation}
  p = 101325
  \left(1 - 2.25577\times 10^{-5}\, h\right)^{5.25588}
  \quad [\mathrm{Pa}].
\end{equation}

Density is the ideal-gas law at that pressure:

\begin{equation}
  \rho = \frac{p}{R T}, \qquad R = 287.05\,\mathrm{J/(kg\,K)}.
\end{equation}

Two densities are frozen for the mission (AeroSandbox ISA at field $+$ 400\,ft AGL, with the design-day temperature offset):

\begin{itemize}
  \item Night: $T = 26^\circ\mathrm{C}$ ($T_\mathrm{amb,min}+2$) $\Rightarrow$ $\rho_N = 1.048\,\mathrm{kg/m^3}$.
  \item Day: $T = 36^\circ\mathrm{C}$ ($T_\mathrm{amb,max}-4$) $\Rightarrow$ $\rho_D = 1.014\,\mathrm{kg/m^3}$.
\end{itemize}

Air viscosity is Sutherland's law at the night temperature, $\mu = 1.84\times 10^{-5}\,\mathrm{Pa\,s}$. Reynolds number of a length $\ell$ is

\begin{equation}
  \mathrm{Re}_\ell = \frac{\rho V \ell}{\mu}.
\end{equation}

\begin{intuition}
\textbf{Intuition.} Hot air is thinner, and El Mirage at $\sim$987\,m MSL is already thinner than sea level. A thinner night means the wing needs a higher $C_L$ (or more speed) to hold the same weight, so induced drag rises. Using a hot dawn is deliberately unkind to night watts. Ground effect is not credited: density is at field $+$ 400\,ft AGL, not a cushion of extra lift near the lake bed.
\end{intuition}

% =====================================================================
\section{Geometry}
% =====================================================================

\subsection{Wing planform}

The wing is three pieces. Inboard, between the booms, the chord is $c_r$. Outboard of each boom the chord stays $c_r$ until a taper-start station, then linearly tapers to $c_\mathrm{tip} = \lambda c_r$. Leading edge unswept, so taper is trailing-edge cut.

Let $y_b$ be half the boom spacing and $L_\mathrm{out} = b/2 - y_b$ the outboard panel length. With taper-start fraction $f$ (0 = start at the boom),

\begin{align}
  L_\mathrm{rect} &= f\, L_\mathrm{out}, \\
  L_\mathrm{tap}  &= (1-f)\, L_\mathrm{out}, \\
  S &= 2 y_b c_r + 2 L_\mathrm{rect} c_r + L_\mathrm{tap}(c_r + c_\mathrm{tip}).
\end{align}

Mean aerodynamic chord for an unswept piecewise-linear planform:

\begin{equation}
  \bar{c} = \frac{2}{S}\int_{-b/2}^{b/2} c(y)^2\,\mathrm{d}y.
\end{equation}

Aspect ratio is the usual $ \mathrm{AR} = b^2 / S $.

Boom spacing and area depend on each other (H-stab span = boom spacing, and $V_H$ uses $S$ and $\bar{c}$), so the code iterates that loop 16 times.

\subsection{Horizontal tail}

Target tail volume

\begin{equation}
  V_H = \frac{S_H \ell_t}{S \bar{c}} = 0.50 \quad \text{\flag{}}.
\end{equation}

When boom spacing is forced (e.g.\ $1.30\,\mathrm{m}$), H-stab span equals that spacing and chord is

\begin{equation}
  c_H = \max\left(\frac{V_H S \bar{c}}{\ell_t\, b_\mathrm{boom}},\;
  c_{H,\min}\right),
\end{equation}

where $c_{H,\min}$ is the chord that still fits one $130\,\mathrm{mm}$ cell row after the elevator keep-out. That is why a short tail arm does not shrink the H-stab to nothing: solar packing floors the chord.

\subsection{Vertical tails}

The V-stab aerodynamic centre sits on its own arm $\ell_v$, independent of the H-stab arm $\ell_t$. Total area (both fins)

\begin{equation}
  V_V = \frac{S_V \ell_v}{S b} = 0.030 \quad \text{(user; both fins together)}.
\end{equation}

There is \textbf{no} fin-area floor at a 1.30\,m arm: stretching $\ell_v$ shrinks $S_V$ at fixed $V_V$. Each fin has area $S_V/2$ and $\mathrm{AR}_V = 1.5$ (held at the start of \code{CONTINUOUS}). No rudder; yaw control is differential thrust. No $\pi$-tail endplate credit on $C_{n_\beta}$.

The V-stab leading edge is required to sit at or behind the H-stab trailing edge:

\begin{equation}
  \ell_v \ge \ell_t + 0.75\, c_H + 0.25\, c_V(\ell_v).
\end{equation}

$c_V$ shrinks as $\ell_v$ grows, so the code iterates to the fixed point and floors any requested $\ell_v$ that would overlap the H-stab in $x$.

\subsection{Boom length}

The carbon tube runs from the wing quarter-chord (at the boom) to the V-stab trailing edge. The fuselage pod covers prop disk $\to$ wing trailing edge and is \emph{not} part of $L_\mathrm{boom}$.

\begin{equation}
  L_\mathrm{boom} = x_{\mathrm{V,TE}} - \tfrac14 c_r,
\end{equation}

where $x_{\mathrm{V,TE}} = \tfrac14 \bar{c} + \ell_v + \tfrac34 c_V$. The H-stab camber is in the \textbf{wing plane} ($z_H = 0$); the fins stand up from the boom further aft. Fuselage length and CAD-scaled wetted area (one pod):

\begin{align}
  L_\mathrm{fuse} &= 0.70\,\mathrm{m} + c_r, \\
  S_{\mathrm{wet,fuse}} &= 173128.421\,\mathrm{mm}^2 \times \frac{L_\mathrm{fuse}}{980\,\mathrm{mm}}.
\end{align}

The $980\,\mathrm{mm} / 173128.421\,\mathrm{mm}^2$ point is from CAD. Linear scale assumes constant section (motor/spinner girth does not grow with chord).

% =====================================================================
\section{Mass}
% =====================================================================

Wetted-area factor $k_w = 2.05$ (both sides plus a little thickness). Wing structure uses $0.84\times 0.8 = 0.672\,\mathrm{kg/m^2}$ of wetted area (user, 17 August 2026: 0.8$\times$ the previous main-wing model only), \emph{including} encapsulation, \emph{excluding} cells. Tails stay at $0.44\,\mathrm{kg/m^2}$ wetted.

\begin{align}
  m_\mathrm{wing}  &= 2.05\, S \times 0.672, \\
  m_{H}            &= 2.05\, S_H \times 0.44, \\
  m_{V}            &= 2.05\, S_V \times 0.44, \\
  m_\mathrm{booms} &= 2 \times L_\mathrm{boom} \times 0.174\,\mathrm{kg/m}, \\
  m_\mathrm{cells} &= N_\mathrm{cells} (0.0065 + 0.0015)\,\mathrm{kg}, \\
  m_\mathrm{MPPT}  &= N_\mathrm{strings} \times 0.108\,\mathrm{kg}, \\
  m_\mathrm{batt}  &= 2 \times 1.278\,\mathrm{kg}.
\end{align}

Boom linear density $0.174\,\mathrm{kg/m}$ is the 1.0\,in OD carbon tube (was $0.290\,\mathrm{kg/m}$ at $\varnothing 38.1\,\mathrm{mm}$). Parasite drag uses the 25.4\,mm OD.

Fixed masses (prototype table): avionics 0.30, power boards 0.10, servos 0.08, ESCs 0.06, wiring 0.096, fuselage pods 0.20, superstructures 0.20, boom/V-stab interfaces 0.10, boom/H-stab fittings 0.06 (same lump as the old $\pi$-tip joint). Total fixed $= 1.196\,\mathrm{kg}$.

Motors and props are catalog masses (two of each).

\begin{examplebox}
\textbf{Worked example --- current nearest-miss} ($b=5.91$, $c_r=0.314$, $\lambda=0.60$, $\ell_t=1.28$, $\ell_v=1.70$, 76 cells).

\begin{center}
\begin{tabular}{@{}lr@{}}
\toprule
Item & Mass [kg] \\
\midrule
Wing structure & 2.25 \\
H-stab & 0.22 \\
V-stabs & 0.15 \\
Booms ($2\times 1.87\,\mathrm{m}\times 0.174$) & 0.65 \\
Cells + interconnect & 0.61 \\
Four MPPTs & 0.43 \\
Two packs & 2.56 \\
Two A40-12L & 0.55 \\
Two APC $16\times12$E & 0.12 \\
Fixed & 1.20 \\
\midrule
\textbf{Total} & \textbf{8.75} \\
\bottomrule
\end{tabular}
\end{center}

The 0.8$\times$ wing areal cut $\sim$0.6\,kg off the previous 0.84\,kg/m$^2$ wing. The 1\,in tube is both lighter per metre and thinner in parasite drag than the 1.5\,in / 0.290\,kg/m boom. The V-stabs are smaller than a co-located $\pi$-tail at $\ell_t$ because $\ell_v$ is longer at fixed $V_V$. Batteries are still the heaviest single lump.
\end{examplebox}

\begin{intuition}
\textbf{Intuition.} You do not ``buy'' span for free. A 6\,m wing at 0.672\,kg/m$^2$ wetted is already $\sim$2.3\,kg before cells. Lengthening $\ell_v$ (not just $\ell_t$) adds carbon on \emph{both} booms, because the tube still runs to the V-stab trailing edge. That extra mass flies all night.
\end{intuition}

% =====================================================================
\section{Aerodynamics}
% =====================================================================

File: \code{solar\_uav/aircraft.py}. Airfoil polars: NeuralFoil \code{xlarge} via AeroSandbox, cached in \code{data/polars/}. Wing section is S4310; tails are NACA~0010. Polar grid: $\mathrm{Re} \in \{50,100,150,200,300\}\times 10^3$, $\alpha \in [-6^\circ, 16^\circ]$ at $0.25^\circ$.

NeuralFoil polars are not monotonic in $C_L$ at negative $\alpha$. The code takes the \textbf{increasing-$C_L$ pre-stall branch} (from $C_{L,\min}$ to $C_{L,\max}$), then interpolates $c_d$, $\alpha$, and $c_m$ versus $C_L$ and $\log\mathrm{Re}$.

\subsection{Level-flight lift}

Dynamic pressure and net lift coefficient (weight on the wing reference area):

\begin{equation}
  q = \tfrac12 \rho V^2, \qquad
  C_{L,\mathrm{net}} = \frac{mg}{q S}.
\end{equation}

The H-stab carries a (usually negative) load, so the wing must make a little more than $C_{L,\mathrm{net}}$:

\begin{equation}
  C_{L,w}\, S + C_{L,H}\, S_H = mg
  \quad \Rightarrow \quad
  C_{L,w} = C_{L,\mathrm{net}} - \frac{S_H}{S} C_{L,H}.
\end{equation}

$C_{L,H}$ is not a guess. It comes from the trim iteration below.

\begin{intuition}
\textbf{Intuition.} $C_{L,\mathrm{net}}$ is ``how hard is this wing working, as a whole, to hold the airplane up?'' If the tail is pushing down, the wing has to make \emph{more} lift than the weight, and induced drag is paid on that higher wing $C_L$. That is why a badly trimmed tail is not free.
\end{intuition}

\subsection{Lift-curve slopes (Helmbold)}

Finite-wing lift slope, per radian:

\begin{equation}
  a = \frac{2\pi}{1 + 2/\mathrm{AR}}.
\label{eq:helmbold}
\end{equation}

Used for the wing ($a_w$) and the H-stab ($a_t$ with the H-stab's actual AR). This is the classical correction that a high-AR wing is closer to $2\pi$ and a stubby tail is much less.

Downwash derivative at the tail, handbook:

\begin{equation}
  \frac{\mathrm{d}\varepsilon}{\mathrm{d}\alpha} = \frac{2 a_w}{\pi\,\mathrm{AR}}.
\end{equation}

\subsection{Oswald efficiency (induced-drag factor)}

Rectangular calibration $e_0 = 0.90$ \flag{}. Taper scales that using AeroSandbox's Nita--Scholz estimate:

\begin{equation}
  e = \min\left(0.98,\;
      e_0 \frac{e_\mathrm{Nita}(\lambda_\mathrm{eq},\mathrm{AR})}
               {e_\mathrm{Nita}(1,\mathrm{AR})}\right).
\end{equation}

$\lambda_\mathrm{eq}$ is an equivalent taper from the actual area and tip chord, clipped to $[0.2, 1]$. The 0.98 cap is a hard ceiling: we do not claim a better-than-nearly-elliptical wing.

On the current $\lambda = 0.60$, $\mathrm{AR} = 21.3$ airplane, this \textbf{hits the 0.98 cap}.

\begin{intuition}
\textbf{Intuition.} Induced drag is the price of making lift with a finite wing: the tip vortices tilt the lift vector backward. Elliptical spanload pays the minimum price, $C_{D_i} = C_L^2/(\pi\,\mathrm{AR})$. Oswald $e$ is ``how close are we to that ellipse?'' $e=0.90$ is an honest rectangle. $e=0.98$ is ``this taper already did most of the job.'' That is why adding washout on top of $\lambda=0.4$ does not buy much --- there is almost no induced-drag left to save, and the cap forbids claiming more.
\end{intuition}

\subsection{Induced drag of the wing}

\begin{equation}
  C_{D_i,w} = \frac{C_{L,w}^2}{\pi\,\mathrm{AR}\, e}.
\label{eq:cdi}
\end{equation}

If washout is zero, that is the whole induced model. If washout is nonzero, $e$ is replaced by an effective $e$ from lifting-line (see \S\ref{sec:washout}).

\subsection{Wing profile drag}

At the local $(C_{L,w}, \mathrm{Re}_{\bar{c}})$, read $c_d$ from the S4310 NeuralFoil polar:

\begin{equation}
  C_{D_{p,w}} = c_{d,\mathrm{S4310}}\!\left(C_{L,w},\, \mathrm{Re}_{\bar{c}}\right).
\end{equation}

This is a 2D section drag treated as constant across the span when washout is zero. With washout, a strip integral adds the difference vs.\ the untwisted wing (again \S\ref{sec:washout}).

\subsection{Horizontal-stabilizer drag}

Profile from NACA~0010 at the tail's $C_{L,H}$ and $\mathrm{Re}_{c_H}$. Induced with a conservative tail Oswald $e_H = 0.70$:

\begin{equation}
  C_{D_{p,H}} = c_{d,0010}(C_{L,H}, \mathrm{Re}_{c_H}), \qquad
  C_{D_{i,H}} = \frac{C_{L,H}^2}{\pi\,\mathrm{AR}_H \cdot 0.70}.
\end{equation}

These coefficients are referenced to $S_H$, not $S$.

\subsection{Vertical-stabilizer drag}

Both fins at $C_L = 0$ (no sideslip in the cruise model):

\begin{equation}
  C_{D_{p,V}} = c_{d,0010}(0, \mathrm{Re}_{c_H}),
\end{equation}

referenced to $S_V$ (the H-stab chord Reynolds number is used as a stand-in for the fin). No induced term: they are not lifting in this model.

\subsection{Parasite drag of pods and booms}

Turbulent skin friction (Blasius-like power law, floored at $\mathrm{Re}=10^4$):

\begin{equation}
  C_f = \frac{0.074}{\mathrm{Re}^{0.2}}.
\end{equation}

Two fuselage pods. Wetted area is interpolated from CAD ($173128.421\,\mathrm{mm}^2$ at $L=980\,\mathrm{mm}$), Reynolds number on the actual pod length $L_\mathrm{fuse} = 0.70\,\mathrm{m}+c_r$, form factor 1.25:

\begin{equation}
  f_\mathrm{pods} = 2\, C_f(\mathrm{Re}_{L_\mathrm{fuse}}) \cdot 1.25 \cdot S_{\mathrm{wet,fuse}}(L_\mathrm{fuse}).
\end{equation}

Two circular booms, wetted area $\pi D L$ per tube, $D = 1.0\,\mathrm{in} = 25.4\,\mathrm{mm}$, form/interference 1.10 on the booms, then another 1.10 on the sum:

\begin{align}
  f_\mathrm{booms} &= 2\, C_f(\mathrm{Re}_L) \cdot 1.10 \cdot (\pi D L_\mathrm{boom}), \\
  f_\mathrm{parasite} &= 1.10\,(f_\mathrm{pods} + f_\mathrm{booms}).
\end{align}

$f$ is a drag area: $D = q f$. No laminar-run credit, no T-tail wake shielding.

\subsection{Putting drag together}

\begin{equation}
\begin{aligned}
  D = q \Big[&
      S \bigl(C_{D_{p,w}} + C_{D_{i,w}}\bigr)
    + S_H \bigl(C_{D_{p,H}} + C_{D_{i,H}}\bigr) \\
    &+ S_V C_{D_{p,V}}
    + f_\mathrm{parasite}
  \Bigr] \times k_\mathrm{trim}.
\end{aligned}
\label{eq:drag}
\end{equation}

$k_\mathrm{trim} = 1.00$. There is no lumped ``trim drag fudge.'' Trim is modeled by actually loading the tail.

\subsection{Incidence and trim}
\label{sec:trim}

The wing is jigged so the \textbf{booms fly at $\alpha \approx 0$} at night min-airspeed. That is a structural / packaging choice (the carbon tubes want to be at zero), not an aerodynamic optimum.

NeuralFoil gives the 2D geometric angle $\alpha_{2D}$ that produces the cruise $C_L$ on S4310. The wing also has an induced angle $C_L/(\pi\,\mathrm{AR}\,e)$. Boom angle of attack is

\begin{equation}
  \alpha_\mathrm{boom} = \alpha_{2D} + \frac{C_L}{\pi\,\mathrm{AR}\,e} - i_w.
\end{equation}

Setting $\alpha_\mathrm{boom}=0$ at the design point:

\begin{equation}
  i_w = \alpha_{2D}(C_L) + \frac{C_L}{\pi\,\mathrm{AR}\,e}.
\label{eq:iw}
\end{equation}

On the current airplane, $i_w \approx 3.70^\circ$.

The H-stab is jigged so the \textbf{elevator is $\approx 0$} at that same point. Handbook downwash at the tail:

\begin{equation}
  \varepsilon = \frac{2 C_{L,w}}{\pi\,\mathrm{AR}}
  \qquad \text{(radians; no T-tail height credit)}.
\label{eq:eps}
\end{equation}

Section $c_m$ about $c/4$ from NeuralFoil, tail volume $V_H$, tail lift slope $a_t$:

\begin{equation}
  i_t = \varepsilon + \frac{c_m}{V_H a_t}
  \qquad \text{(angles in consistent units)}.
\label{eq:it}
\end{equation}

On the current airplane, $i_t \approx 0.07^\circ$.

Inside \code{drag\_n}, three iterations:

\begin{enumerate}
  \item From current $C_{L,w}$, get $\alpha_{2D}$ and $\alpha_i = C_{L,w}/(\pi\,\mathrm{AR}\,e)$.
  \item Boom $\alpha_f = \alpha_{2D} + \alpha_i - i_w$.
  \item $C_{L,H} = a_t\,(\alpha_f + i_t - \varepsilon)$.
  \item $C_{L,w} \leftarrow C_{L,\mathrm{net}} - (S_H/S) C_{L,H}$.
\end{enumerate}

\begin{intuition}
\textbf{Intuition for \eqref{eq:eps}.} A lifting wing throws air down. The tail flies in that descending wake, so it sees a more negative angle than the boom. $\varepsilon \approx 2C_L/(\pi\,\mathrm{AR})$ is the standard horseshoe-vortex estimate: more lift or less AR $\Rightarrow$ stronger downwash. A T-tail sits higher, in weaker downwash. The H-stab is in the \textbf{wing plane}, so there is no height credit to take. If VLM at that $z_H=0$ is \emph{stronger} than handbook, production uses VLM (currently $k_\varepsilon \approx 1.21$). We never keep a $k_\varepsilon<1$ ``T-tail'' gift.
\end{intuition}

AeroSandbox is used as a \textbf{conservative envelope}, not as the scoring polar (\code{solar\_uav/asb\_physics.py}):

\begin{itemize}
  \item Atmosphere is the AeroSandbox ISA at field $+$ 400\,ft AGL, with the design-day temperature offset.
  \item Wing-only VLM downwash at the real H-stab height ($z_H=0$): if it is \emph{weaker} than \eqref{eq:eps}, keep the handbook (no T-tail credit). If it is stronger, use VLM. Induced $e$ is $\min(e_\mathrm{handbook}, e_\mathrm{VLM})$.
  \item AeroBuildup $C_{l_p}$, $C_{n_\beta}$, $C_{n_r}$: keep the number that makes roll rate, weathercock, or yaw damping \emph{harder}.
  \item Pitch inertia is $\max($lumped, CAD-located \code{MassProperties} $I_{yy})$. Area-painted thin-shell inertia is printed as a FLAG and not used.
  \item XFoil vs NeuralFoil at cruise $C_L$: inflate section $c_d$ if XFoil is draggier; never deflate. If the XFoil binary is missing, FLAG and keep NeuralFoil.
\end{itemize}

Run \code{scripts/run\_asb\_audit.py} to print the envelope on the reference airplane.

\subsection{Static margin}

CG is locked at the wing quarter-chord \flag{}. Neutral-point shift is then entirely the tail:

\begin{equation}
  \mathrm{SM}
    = V_H \frac{a_t}{a_w}
      \left(1 - \frac{\mathrm{d}\varepsilon}{\mathrm{d}\alpha}\right).
\end{equation}

Target was 12\% MAC; minimum accepted is 8\%. The current airplane sits at \textbf{43\% MAC} --- very stable. $V_H$ comes out $0.69$ rather than the $0.50$ volume target because the H-stab chord is floored by one $130\,\mathrm{mm}$ cell row after the elevator keep-out. That extra tail is packing, not an SM request.

Elevator effectiveness is the empirical stand-in $\tau = \sqrt{c_e/c}$ \flag{}. Elevator must trim a $\Delta C_L = 1$ at $\le 15^\circ$. Pitch acceleration about the CG must reach $0.30\,\mathrm{rad/s^2}$ ($\sim 20^\circ$ in 1.5\,s from rest) at min airspeed.

\subsection{Stall and the speeds we actually fly}

2D $C_{L,\max}$ is the peak of the NeuralFoil pre-stall branch. 3D knockdown:

\begin{equation}
  C_{L,\max,3\mathrm{D}} = 0.90\, C_{L,\max,2\mathrm{D}}.
\end{equation}

For S4310, $C_{L,\max,2\mathrm{D}} \approx 1.44$, so $C_{L,\max,3\mathrm{D}} \approx 1.30$. Stall speed:

\begin{equation}
  V_s = \sqrt{\frac{2mg}{\rho S C_{L,\max,3\mathrm{D}}}}.
\end{equation}

Operational floor (flown loiter speed):

\begin{equation}
  V_{\min} = \max\bigl(1.3\, V_s,\; 7.7\,\mathrm{m/s}\bigr).
\end{equation}

$1.2\,V_s$ remains a never-below stall margin. The 89-hour flight is a banked loiter, so the energy march pays $1.3\,V_s$. The $7.7\,\mathrm{m/s}$ is 15\,kt of required crosswind capability. On the current airplane $1.3 V_s = 11.42\,\mathrm{m/s}$ and \textbf{binds}: unconstrained aero min-power is $\sim$9.2\,m/s ($DV \approx 29\,\mathrm{W}$ vs.\ $36\,\mathrm{W}$ at the floor). That extra 2.2\,m/s is about $+10\,\mathrm{W}$ of night bus.

Cruise is \textbf{not} $V_{\min}$ blindly. It is the speed $\ge V_{\min}$ that minimises shaft power $P_\mathrm{aero} = D V$. The 1-D search is bounded to $8$--$18\,\mathrm{m/s}$, starts at $12\,\mathrm{m/s}$, and refines to $0.01\,\mathrm{m/s}$ (Brent), still floored by the loiter margin. On the current airplane the floor \textbf{is} the min-power point: unconstrained aero min-power is $\sim$9.2\,m/s, but $1.3 V_s = 11.42\,\mathrm{m/s}$ binds.

\begin{intuition}
\textbf{Why min power, not max $L/D$?} Endurance on a battery cares about \emph{watts}, not distance. Max $L/D$ minimises drag (energy per metre). Min $D V$ minimises energy per second. For a solar loiterer the clock is the enemy, so we fly min power. The extra $1.3$ factor (vs.\ $1.2$) is so a $20^\circ$ bank does not eat the stall margin --- $\sqrt{1/\cos 20^\circ}\approx 1.03$, so $1.3$ already covers turn-stall.
\end{intuition}

Climb is the ideal excess-power expression at day min-power speed, using max thrust available on 25.2\,V:

\begin{equation}
  \dot{h} = V \frac{T - D}{mg}.
\end{equation}

No acceleration, no prop-in-climb polar change, no takeoff ground roll. On the 12L + APC $16\times12$E this is $\sim 15\,\mathrm{m/s}$ versus a $1.5\,\mathrm{m/s}$ requirement. Climb is not the constraint. It only vetoes fans that are efficient at cruise but cannot pull.

\subsection{Ailerons, elevons, and fins}
\label{sec:control}

The performance polar does not fly the airplane. Roll, weathercock, and yaw damping are hard constraints so a later sim cannot inherit a 5$^\circ$/s placeholder.

\textbf{Ailerons} are 25\% local chord. On a taper they start where the second cell row drops (or further inboard if 35$^\circ$/s is not met) and carry no cells on the moving surface. On a rectangle, cells may cover the aileron. \textbf{Elevons} are the existing elevator, split left/right, used differentially for extra roll. There is \textbf{no rudder}; yaw is differential thrust.

Steady roll helix (full $20^\circ$ aileron + $20^\circ$ elevon):

\begin{equation}
  \frac{p b}{2V} = -\frac{C_{l_{\delta a}}\delta_a + C_{l_{\delta e}}\delta_e}{C_{l_p}},
  \qquad
  C_{l_{\delta a}} = \frac{2 a_w \tau}{S b}\int_{y_1}^{y_2} c(y)\, y\,\mathrm{d}y.
\end{equation}

$C_{l_p}$ is a high-AR handbook estimate inflated 15\% so the rate is harder to achieve. Acceptance: $p \ge 35^\circ/\mathrm{s}$ at loiter speed ($\sim 20^\circ$ bank in about a second). A 10\,cm tip strip on this wing makes $\sim 5^\circ/\mathrm{s}$ --- that was the sim.

Weathercock and yaw damping (two NACA~0010 fins, Helmbold $a_v$, $\eta=0.90$, no endplate credit):

\begin{equation}
  C_{n_\beta} = \eta a_v V_V, \qquad
  C_{n_r} = -2\eta a_v \frac{S_V}{S}\left(\frac{\ell_v}{b}\right)^2.
\end{equation}

Floors: $C_{n_\beta}\ge 0.10$, $|C_{n_r}|\ge 0.030$ are \textbf{off} (differential thrust is the yaw effector). Derivatives are still computed. Hinge/side-gap drag of ailerons and elevons is added to the polar.

Differential-thrust yaw (one motor $T_{\max}$, the other idle) must overpower weathercock at $5^\circ$ $\beta$. The \textbf{corrective rate} is the still-air heading rate with loiter-preserving $\Delta T$ (one motor up, the other down, sum held at drag) versus $C_{n_r}$:

\begin{equation}
  r = \frac{N_{\mathrm{DT,loiter}}}{q S b\, |C_{n_r}|\, b/(2V)},
  \qquad
  N_{\mathrm{DT,loiter}} = 2\min(T_{\mathrm{cruise}},\, T_{\max}-T_{\mathrm{cruise}})\, y_{\mathrm{boom}}.
\end{equation}

Acceptance: $r \ge 5^\circ/\mathrm{s}$ (unwind $5^\circ$ in about a second). $C_{n_r}$ for this rate is the \emph{more negative} of handbook vs AeroBuildup. $I_{zz}$ sizes the initial $\dot r$ and the 1-DOF time to yaw $5^\circ$; motor spool-up is not modelled, so those two are optimistic. Max-effort (idle one motor) rate is reported and is not the gate.

\textbf{No dihedral} is modelled. Spiral-mode static stability is not claimed; the autopilot is assumed to hold a slow spiral with the sized ailerons.

\subsection{Washout}
\label{sec:washout}

Geometric washout is tip leading-edge down, starting at a station on the outboard panel. \textbf{Zero washout does not change the lumped polar at all} --- the lifting-line code is not even called. That is deliberate, so the published baseline cannot move when someone toggles a unused twist flag.

When twist is nonzero, a Fourier lifting-line (Prandtl / Anderson, 11 odd modes on the half-span, file \code{solar\_uav/lifting\_line.py}) solves

\begin{equation}
  \Gamma(\theta) = 2 b V \sum_{n=1,3,\ldots} A_n \sin(n\theta),
  \qquad y = (b/2)\cos\theta.
\end{equation}

Section lift $c_\ell = 2\Gamma /(V c)$. Wing $C_L = \pi\,\mathrm{AR}\, A_1$. Trefftz-plane Oswald:

\begin{equation}
  e_\mathrm{LL} = \frac{A_1^2}{\sum n A_n^2}.
\end{equation}

The model does \emph{not} replace handbook $e$ with $e_\mathrm{LL}$ outright (that would let an inviscid method invent a better wing than the calibration). It only reports the \emph{ratio} of twisted to untwisted lifting-line $e$, times the handbook value, still capped at 0.98:

\begin{equation}
  e_\mathrm{eff} = \min\left(0.98,\;
    e_\mathrm{handbook} \frac{e_\mathrm{LL,twisted}}{e_\mathrm{LL,untwisted}}\right).
\end{equation}

Strip profile drag is integrated at the collocation stations; only the \emph{increment} vs.\ the same planform untwisted is added. Stall becomes first-strip $C_{L,\max}$, and cannot exceed the untwisted 3D cap.

\begin{examplebox}
\textbf{Example --- why the search refused washout.}
On $\lambda = 0.4$ (and on the current $\lambda \approx 0.60$) the spanload is already near-elliptical, $e$ is already at the 0.98 cap. Twist cannot raise $e$ further. It \emph{can} unload the tip (good for stall) but the first-strip rule and the profile-drag increment make cruise slightly worse. The current search parked washout at $\approx 0^\circ$. On a \emph{rectangular} wing, washout does raise $e$ ($0.90 \to \sim 0.95$), but flying a rectangle to get that twist benefit also adds tip chord, cells, mass, and drops AR. That package still lost on margin. The optimizer's message: \textbf{taper already bought the induced-drag win; do not pay mass for a second copy of the same win.}
\end{examplebox}

\subsection{Worked drag budget (current night cruise)}
\label{sec:dragbudget}

$m = 8.745\,\mathrm{kg}$, $V = 11.42\,\mathrm{m/s}$ ($1.3 V_s$ binds), $\rho = 1.048$, $q = 68.3\,\mathrm{Pa}$, $S = 1.636\,\mathrm{m^2}$.

\begin{center}
\begin{tabular}{@{}lrrr@{}}
\toprule
 & $C_L$ / $C_D$ & Drag [N] & Share \\
\midrule
Net airplane $C_L$ & 0.768 & & \\
Wing $C_L$ & 0.784 & & \\
H-stab $C_L$ & $-0.109$ & & \\
Wing profile (S4310, $\mathrm{Re}_{\bar{c}}\approx 1.84\times 10^5$) & $C_{D_p}=0.0123$ & 1.38 & 44\% \\
Wing induced ($e=0.98$) & $C_{D_i}=0.0094$ & 1.05 & 33\% \\
H-stab profile & & 0.20 & 6\% \\
H-stab induced ($e_H=0.70$) & & 0.018 & 1\% \\
V-stab profile ($C_L=0$) & & 0.14 & 4\% \\
Pods & & 0.17 & 5\% \\
Booms ($\varnothing 25\,\mathrm{mm}$, $L=1.87\,\mathrm{m}$) & & 0.11 & 4\% \\
Control-surface gaps & & 0.10 & 3\% \\
\midrule
\textbf{Total} & $C_D=0.0283$ on $S$ & \textbf{3.16} & \\
Aero power $DV$ & & \textbf{36.1 W} & \\
$L/D$ & 27.1 & & \\
\bottomrule
\end{tabular}
\end{center}

\begin{intuition}
\textbf{Read this table like a designer, not an aerodynamicist showing off.} About 77\% of cruise drag is the wing. Profile now beats induced (44\% vs.\ 33\%) because the airplane is lighter and is forced to fly $1.3 V_s$ instead of aero min-power. Unconstrained min-power is $\sim$9.2\,m/s and $DV \approx 29\,\mathrm{W}$; the speed floor costs about $7\,\mathrm{W}$ of aero and $\sim$10\,W of bus. The entire empennage plus pods plus booms plus gaps is $\sim$0.74\,N. Polishing boom fairings will not close a 158\,Wh night gap. The remaining polar levers at fixed mass are $C_{L,\max}$ (or relaxing $1.3\to 1.2$) and wing profile $C_D$. Span and $e$ are already maxed.
\end{intuition}

Aero power 36\,W is \emph{not} night bus power. The propeller, motor, ESC, and the 5\,W avionics sit between $DV$ and the battery. That chain currently turns 36\,W of aero into $70.2\,\mathrm{W}$ of propulsion plus $5\,\mathrm{W}$ avionics $= 75.2\,\mathrm{W}$ on the bus. Propulsion, not the polar, is where most of the remaining watts live.

% =====================================================================
\section{Solar array}
% =====================================================================

Files: \code{solar\_uav/environment.py}, \code{solar\_uav/components/solar\_array.py}.

\subsection{Sun}

pvlib Ineichen clear-sky at the site, minute resolution. The array is \textbf{horizontal} (wing top). Angle of incidence equals the solar zenith angle.

Plane-of-array irradiance before cells:

\begin{equation}
  E_\mathrm{POA}
    = \underbrace{E_\mathrm{DNI}\cos\zeta \cdot \mathrm{IAM}(\zeta)}_{\text{beam, with reflection}}
    + \underbrace{0.95\, E_\mathrm{DHI}}_{\text{diffuse IAM \flag{}}}.
\end{equation}

Then encapsulation transmission 0.97 \flag{}:

\begin{equation}
  E_\mathrm{eff} = 0.97\, E_\mathrm{POA}.
\end{equation}

IAM is pvlib's physical air/glass model. A horizontal wing at El Mirage ($34.65^\circ\mathrm{N}$) never sees the sun overhead except near noon in June, and even then $\zeta \ne 0$ most of the day. That cosine is the main reason ``battery night'' is longer than civil night.

\subsection{Cell temperature}

No flight-cooling credit. NOCT-style

\begin{equation}
  T_\mathrm{cell} = T_\mathrm{amb} + k\, E_\mathrm{eff},
\end{equation}

with $k$ calibrated so solstice-noon static $T_\mathrm{cell} = 75^\circ\mathrm{C}$ (user requirement). Ambient is a sinusoid, min at 05:00, max at 15:30, between 24 and $40^\circ\mathrm{C}$.

\subsection{Electrical output}

C60 bin I \flag{}: $P_\mathrm{mp} = 3.40\,\mathrm{W}$ at STC. Temperature coefficient $\gamma = -0.32\%/^\circ\mathrm{C}$. Wiring / mismatch / soiling 0.965 \flag{}. MPPT 0.95.

\begin{equation}
  P_\mathrm{bus}
    = N_\mathrm{cells}\, P_\mathrm{mp}
      \frac{E_\mathrm{eff}}{1000}
      \bigl[1 + \gamma (T_\mathrm{cell}-25)\bigr]
      \times 0.965 \times 0.95.
\end{equation}

\subsection{Why 26 cells is a hard wall}

Boost MPPT: string $V_\mathrm{oc}$ must stay below the depleted 6S bus (19.2\,V) so the converter can always boost. Cold cells ($10^\circ\mathrm{C}$) have the highest Voc. Bin I: $V_\mathrm{oc}(10^\circ\mathrm{C}) \approx 0.713\,\mathrm{V}$, so

\begin{equation}
  N_{\max} = \left\lfloor \frac{19.2}{0.713} \right\rfloor = 26.
\end{equation}

A rectangular 6.0$\times$0.36\,m outboard wants more than 26 cells in a single series string. Under one-string-per-bay that planform is illegal unless you add strings (extra MPPTs, extra mass) or taper so the tip drops cells.

\begin{examplebox}
\textbf{Example --- extra cells lost.} A mixed-pack rectangular wing with 92--104 cells and extra strings was tried (``just put more solar on''). AR fell to $\sim$16.7, mass rose $\sim$1.4\,kg, night watts rose, and margin got \emph{worse} (about $-184\,\mathrm{Wh}$ vs.\ $-160\,\mathrm{Wh}$ on the tapered 82-cell wing, at the old 410\,g/m booms). Noon already spills. The night does not care that you had 20 more cells at 13:00. It cares that you are heavier and draggier at 02:00.
\end{examplebox}

% =====================================================================
\section{Battery}
% =====================================================================

File: \code{solar\_uav/components/battery.py}.

Two GOLD V1 packs in parallel: 483\,Wh, 23.68\,Ah, 1.278\,kg, 6\,A charge, 313\,W sustained discharge each. Usable energy at the 20\% floor:

\begin{equation}
  E_\mathrm{usable} = 2 \times 483 \times 0.80 = 773\,\mathrm{Wh}.
\end{equation}

SOC is energy-based. Open-circuit voltage is a generic Si-anode Li-ion curve \flag{}, 3.0--4.2\,V/cell, times 6. Internal resistance $45\,\mathrm{m}\Omega$ per pack \flag{}, paralleled.

Terminal power during discharge: $P = V_\mathrm{oc} I - I^2 R$. Chemistry energy removed is $V_\mathrm{oc} I \Delta t$ (you pay $I^2R$ from the bus but it does not come out of SOC as extra capacity --- it is a conversion loss). Charge is the reverse, capped at 6\,A/pack, with a CV taper as $V_\mathrm{oc}$ approaches 25.2\,V.

The march \emph{allows} SOC to fall through 20\% so we can measure how far a miss is. Closure still requires staying above 20\%.

% =====================================================================
\section{The 24-hour energy march}
% =====================================================================

File: \code{solar\_uav/mission.py}. Default step: 5 minutes.

Night and day bus loads are each a single number: propulsion at that density's min-power speed, plus 5\,W avionics. Then, at every timestep,

\begin{align}
  \text{is\_night} &= (P_\mathrm{solar} < P_\mathrm{night}), \\
  P_\mathrm{load} &= \begin{cases}
      P_\mathrm{night} & \text{if night}, \\
      P_\mathrm{day} & \text{if day},
    \end{cases} \\
  P_\mathrm{net} &= P_\mathrm{solar} - P_\mathrm{load}.
\end{align}

\begin{warningbox}
\textbf{Do not mentally replace $P_\mathrm{net}$ with $-P_\mathrm{night}$ at dusk.} If the array is making 60\,W and the load is 75\,W, the pack only supplies 15\,W. Calling that ``night'' and draining 75\,W is how people convince themselves the physics is worse than it is --- and then over-correct with a third pack. The code already counts the shoulders.
\end{warningbox}

The march starts at the \textbf{last afternoon timestep} where $P_\mathrm{net} > 1\,\mathrm{W}$, with SOC $= 1$, then wraps 24 hours. That way the night is fully counted and we test refill on the following afternoon.

\textbf{Battery-true night} is hours with $P_\mathrm{solar} < P_\mathrm{night}$. On 21 June at El Mirage ($34.65^\circ\mathrm{N}$) that is \textbf{14.0 hours} on the current airplane, not the $\sim$9.5 hours of civil night. A horizontal wing at this latitude simply does not make cruise power at low sun.

\begin{examplebox}
\textbf{Example --- why ``75\,W $\times$ 14\,h = 1050\,Wh, we need another pack'' is the wrong arithmetic.}
Usable pack is 773\,Wh. Shoulders (dawn/dusk when the array makes some but not all of cruise) contribute on the order of 100+\,Wh that a naive $P\times t$ ignores. The honest integral is still larger than 773\,Wh --- that is why we miss --- but the scored gap is $158\,\mathrm{Wh}$ below the 20\% floor (and a separate $90\,\mathrm{Wh}$ refill miss), not 300\,Wh. Closing is ``about 11\,W off the night bus'' or ``a bit more pack,'' not ``the polar must be twice as good.''
\end{examplebox}

% =====================================================================
\section{Propulsion}
\label{sec:propulsion}
% =====================================================================

Files: \code{solar\_uav/components/propeller.py}, \code{mejzlik.py}, \code{motor.py}, \code{propulsion.py}.

Two motors, two props, thrust split equally. 6S bus, nominal 22.2\,V, climb uses 25.2\,V.

\subsection{Propeller maps}

APC PER3 files (vortex theory) and Mejzlik 2-blade electric maps (in-house BEM / vortex). Both are \textbf{computational, not measured}. Uncertainty factors, applied in \code{PropulsionSystem} (not inside a \code{nominal=True} table read):

\begin{equation}
  T_\mathrm{used} = 0.92\, T_\mathrm{map}, \qquad
  P_\mathrm{used} = 1.08\, P_\mathrm{map}
  \quad \text{\flag{}}.
\end{equation}

Level flight asks the map for $T_\mathrm{need}/0.92$ so the derated thrust still equals drag.

Nondimensional definitions (McCormick / MIT 16.unified):

\begin{equation}
  J = \frac{V}{n D}, \qquad
  T = C_T(J)\, \rho n^2 D^4, \qquad
  P = C_P(J)\, \rho n^3 D^5, \qquad
  \eta = \frac{J C_T}{C_P}.
\end{equation}

$n$ is rev/s, $D$ is diameter in metres. $C_T$ and $C_P$ come from the published map. Density scaling of a tabulated $T,P$ is $\rho/\rho_\mathrm{ref}$ with $\rho_\mathrm{ref} = 1.225\,\mathrm{kg/m^3}$ for APC.

\textbf{Query rule, in order:}

\begin{enumerate}
  \item If this airspeed exists on an RPM block of the table, interpolate $T(\mathrm{RPM})$ at that speed, density-scale, and pick the RPM that makes the required thrust. \emph{If the slowest tabulated RPM that covers this speed already makes enough thrust, stop. Do not slow down further.}
  \item Otherwise fall back to $C_T(J)$ similarity, but only for $J < J_\mathrm{windmill}$ (the advance ratio where $C_T$ crosses zero). Past windmill the blade is a drag device; there is no cruise solution there.
  \item If table and similarity disagree by more than 15\% (\code{PROP\_SANITY\_REL}), raise a \flag{} on the result. The mission still uses the table value when the table was the source.
\end{enumerate}

\begin{warningbox}
\textbf{The $n^2$-scaling bug, as a worked example.}
At fixed $V$, a naive ``I only need a tiny thrust, scale RPM with $\sqrt{T}$'' walks $J = V/(nD)$ \emph{up}. For the APC $22\times12$E, zero-thrust is near $J \approx 0.67$. Night cruise at 10.7\,m/s wants $n \approx 2070$\,RPM ($J \approx 0.58$) on the table. Scaling down to 1190\,RPM puts $J \approx 1.14$, past windmill. Similarity then says $T=0$. The interpolator had been reporting a few newtons and $\sim$62\,W anyway. That is how a closed design appeared out of nothing. The rule ``never $n^2$-scale RPM down at fixed speed once a covering block exists'' exists so this cannot happen again.
\end{warningbox}

\begin{intuition}
\textbf{Intuition for $J$.} $J$ is ``how much air is rushing at the prop compared to how fast the blades are spinning.'' Small $J$: hovering / climb, blades see mostly rotation, lots of thrust, lots of power. Large $J$: the airplane has run out in front of the blades; they start to windmill. A high-pitch ``cruise prop'' reaches windmill later (larger $J$), which is why we shortlist pitch ratios $0.40$--$0.85$. Even then, a 21$\times$14 at 10.5\,m/s and 1545\,RPM is at $J \approx 0.76$, and Mejzlik similarity says $C_T \approx 0$ while the table still lists $\sim$1.8\,N per blade --- that is why older nearest-misses on MJ-21$\times$14EL were on the optimistic side of the table. The current airplane flies an APC $16\times12$E at 11.42\,m/s and 2365\,RPM; the FLAG there is shaft power, table $28.0\,\mathrm{W}$ vs.\ similarity $22.8\,\mathrm{W}$ (19\%). Night watts use the table. The FLAG is in the CSV for a reason.
\end{intuition}

\subsection{Motor and ESC}

First-order DC motor. Datasheet $K_v$, $R_i$, $I_0$ at 8.4\,V. Torque constant $K_t = 1 / K_{v,\mathrm{rad}}$. Gearbox (if any) ratio $g$, efficiency 0.96 \flag{}. Direct drive pays no gearbox loss.

Given propeller torque $Q_p$ and RPM $n_p$:

\begin{align}
  Q_m &= \frac{Q_p}{g\,\eta_\mathrm{gear}}, \qquad
  n_m = n_p\, g, \\
  I &= \frac{Q_m}{K_t} + I_0, \\
  V &= I R_i + \frac{n_m}{K_v}, \\
  P_\mathrm{bus,1} &= \frac{V I}{\eta_\mathrm{ESC}}.
\end{align}

$I_0$ is still paid when loaded --- it is iron + friction, not ``only at idle.'' Two motors: $P_\mathrm{bus} = 2 P_\mathrm{bus,1}$. A point with $V > 22.2\,\mathrm{V}$ is infeasible at cruise (climb is allowed to 25.2\,V).

$\eta_\mathrm{ESC} = 0.95$ \flag{}. The code does \emph{not} currently inflate that loss when duty is below 50\%, but \code{ESC\_DUTY\_MIN = 0.50} exists as a documented warning. On the current 12L + $16\times12$E at night, $V \approx 6.0\,\mathrm{V}$ at 2365\,RPM, i.e.\ duty $\sim$27\% on a 22.2\,V pack. \textbf{Real ESC efficiency at that duty is worse than 95\%.} Treating 0.95 here is an optimistic FLAG on night bus. A rewind to $K_v \approx 180$--$220$ would put duty at or above 50\%; at fixed $(Q,\mathrm{RPM})$ the 12L rewind model says cruise bus watts are nearly independent of $K_v$, so the catalog kv410 is not a modelled watt win.

\begin{intuition}
\textbf{Intuition for $I = Q/K_t + I_0$.} $Q/K_t$ is the current that actually makes torque. $I_0$ is the current the motor drinks just to spin itself --- hysteresis, windage, bearings. At climb, $Q/K_t$ dwarfs $I_0$ and you barely notice. At night cruise the torque is tiny, so $I_0$ is a large fraction of the bill. That is why a high-$I_0$ motor can look fine on a power-limit datasheet and still lose the night.
\end{intuition}

\subsection{Free-$K_v$ rewind model}

Not a catalog part. Same A40-12L stator, $R_i \propto 1/K_v^2$, $I_0 \propto K_v$. Used to decide what $K_v$ to shop for. At fixed $(Q, n)$, cruise bus power is nearly independent of $K_v$ once you are voltage-feasible --- winding changes trade voltage for current. The 12L wins against the geared 10L because of \textbf{$I_0$ and gearbox loss}, not because 410\,rpm/V is magic.

\subsection{Motors evaluated}
\label{sec:motors}

Catalog in \code{MOTOR\_CATALOG}:

\begin{center}
\begin{tabular}{@{}lrrrrrr@{}}
\toprule
Name & $K_v$ & $R_i$ [$\Omega$] & $I_0$ [A] & $g$ & Mass [kg] & $P_\max$ [W] \\
\midrule
A40-10L V2 kv1100 + 6.7:1 PG & 1100 & 0.019 & 1.70 @ 8.4\,V & 6.7 & 0.315 & 2200 \\
A40-10S V4G kv1600 + 6.7:1 PG & 1600 & 0.012 & 2.75 @ 8.4\,V & 6.7 & 0.247 & 1500 \\
A40-12L V4 kv410 direct & 410 & 0.039 & 1.04 @ 8.4\,V & 1.0 & 0.275 & 1100 \\
\bottomrule
\end{tabular}
\end{center}

Plus the rewind grid $K_v = 180, 200, \ldots, 900$ on the 12L stator (direct drive).

\begin{examplebox}
\textbf{Example --- why 12L beat 10L.}
The 10L is the Hacker ``gliders to 14\,kg, props to 23\,in'' recommendation. It is \emph{not} the search default. At night the prop wants $\sim$1500--2000\,RPM and a few tenths of a newton-metre. After the 6.7:1 gearbox the \emph{motor} is spinning at 10--13\,kRPM and still paying $I_0 = 1.7\,\mathrm{A}$ plus 4\% gearbox. The 12L's $I_0 = 1.04\,\mathrm{A}$ is at the shaft RPM you actually want, with no gearbox. Climb on the 12L + $16\times12$E is still $\sim$15\,m/s. The 10S is worse still ($I_0 = 2.75\,\mathrm{A}$). \textbf{Mindset: size the motor for night $I_0$ and map bite, not for takeoff marketing.}
\end{examplebox}

\subsection{Propellers evaluated}
\label{sec:props}

Shortlist rule (\code{shortlist\_props}):

\begin{itemize}
  \item APC style containing ``E'' (electric), excluding MR, SF, W, N, F2B, 3-blade, 4-blade.
  \item Mejzlik 2-blade electrics (style contains E).
  \item Diameter $\ge 16\,\mathrm{in}$ (no upper cap; climb + cruise decide).
  \item Pitch / diameter $\in [0.40, 0.85]$.
\end{itemize}

That is \textbf{51} maps. Full list in Appendix~\ref{app:props}.

Neighborhood searches used a 7-prop subset that had already survived earlier ranks:

\begin{center}
\texttt{MJ-21x14EL}, \texttt{16x12E}, \texttt{17x12E}, \texttt{18x12E},
\texttt{19x12E}, \texttt{20x13E}, \texttt{21x13E}.
\end{center}

The inner loop currently picks \texttt{16x12E} on the 12L nearest-miss. Larger tabulated diameters cost night watts on-map; similarity-filled points (e.g.\ a $20\times15$E off the table) are excluded. Mejzlik 21-inch fans still appear in the shortlist and in older CSVs, but they are not the current pick.

APC thin-electric mass (when not in a catalog) is the fit $m = 0.263\times 10^{-3}\, D_\mathrm{in}^{1.96}$ \flag{}.

% =====================================================================
\section{Current nearest-miss (does not close)}
% =====================================================================

From \code{outputs/phase4\_candidates.csv} (812 scored, 0 closed), best row of the 17 August 2026 search after independent tails, $0.8\times$ wing areal mass, and 1\,in / $0.174\,\mathrm{kg/m}$ booms. Those numbers are also \code{REF\_*} in \code{config.py} for the visualizer. Washout chose $\approx 0^\circ$.

\begin{center}
\begin{tabular}{@{}ll@{}}
\toprule
Item & Value \\
\midrule
Motor & A40-12L V4 kv410 direct \\
Prop & APC $16\times12$E \\
Wing & $5.907\times 0.314\,\mathrm{m}$, $\lambda=0.598$ from 0.275 of the outboard \\
Area / AR / $e$ & $1.636\,\mathrm{m^2}$, 21.3, 0.98 (cap) \\
H-stab / V-stab arm & $1.278\,\mathrm{m}$ / $1.696\,\mathrm{m}$ (V LE $0.19\,\mathrm{m}$ behind H TE) \\
Boom spacing / elevator & $1.133\,\mathrm{m}$ / 34\% chord \\
$V_H$ actual / $V_V$ & $0.687$ (packing chord floor) / $0.030$ \\
H-stab height / downwash $k$ & $z_H=0$ (wing plane) / $k_\varepsilon=1.21$ (VLM $>$ handbook) \\
Static margin & 43\% MAC \\
Array & 76 cells, strings $26+26+16+8$ (4 MPPTs) \\
Mass & $8.745\,\mathrm{kg}$ (booms $1.868\,\mathrm{m}$, $\varnothing 25.4\,\mathrm{mm}$) \\
Night / day bus & $75.2\,\mathrm{W}$ / $76.5\,\mathrm{W}$ (includes 5\,W avionics) \\
$V_\mathrm{stall}$ / $V_\mathrm{night}$ & $8.78$ / $11.42\,\mathrm{m/s}$ ($1.3 V_s$ binds) \\
Climb & $14.8\,\mathrm{m/s}$ (requirement 1.5) \\
Battery-true night & 14.0\,h \\
$\mathrm{SOC}_\mathrm{min}$ / $\mathrm{SOC}_\mathrm{end}$ & 0.036 / 0.887 \\
Margin & $\mathbf{-158\,\mathrm{Wh}}$ \\
Closed? & No \\
\bottomrule
\end{tabular}
\end{center}

Mass lumps (kg): batteries 2.56, wing 2.25, fixed 1.20, booms 0.65, solar 0.61, motors 0.55, MPPTs 0.43, H-stab 0.22, V-stabs 0.15, props 0.12.

The same search family is lighter than the old 1.5\,in / $0.290\,\mathrm{kg/m}$ $\pi$-tail nearest-miss ($\sim$9.8\,kg, $-112\,\mathrm{Wh}$ at a slower $V_\mathrm{night}$). Independent $\ell_v$, thinner booms, and the $0.8\times$ wing cut mass, but El Mirage density plus the $1.3 V_s$ floor put night bus back at 75\,W and the SOC floor at 0.036. Still about 11\,W of night bus short of the 20\% floor, before you even discount the prop-map and ESC-duty FLAGs.

Prop FLAG on this row: at 2365\,RPM, 11.42\,m/s, table shaft $28.0\,\mathrm{W}$ vs.\ similarity $22.8\,\mathrm{W}$ (19\%). Night watts use the table. ESC duty is $\sim$27\%; the 95\% ESC efficiency is optimistic.

% =====================================================================
\section{Top learnings: design priority and mindset}
% =====================================================================

These are not slogans. They are what the model has already spent compute to tell us.

\subsection{Night is the only constraint that is close}

Day array $\sim$200\,W vs.\ $\sim$75\,W cruise. Climb is an order of magnitude above requirement. MTOW has $\sim$3\,kg of unused budget on the nearest-miss. Span is at the wall. The SOC floor is what fails (refill also misses). \textbf{Every other subsystem is a servant of night bus watts and night duration.}

\subsection{Mass and induced drag beat parasite tricks}

Equation~\eqref{eq:cdi} with $C_L \propto W/S$ means

\begin{equation}
  D_i \propto \frac{W^2}{\rho V^2 b^2 e}.
\end{equation}

Span is maxed. $e$ is capped. At \emph{fixed} mass the remaining polar lever is the $1.3 V_s$ speed floor (true aero min-power is $\sim$9.2\,m/s). If mass is still free, a kilogram of boom is not ``structure we can live with''; it is $C_L$ we fly for 14 hours. Fairings, trip dots, and boom $\varnothing$ are a few percent of drag. Do not spend the team's year there until night bus is honest and the structure is light.

\subsection{The propeller map bite is bigger than the polar}

Aero power is 36\,W. Bus is 75\,W. The missing $\sim$39\,W is prop efficiency at cruise $J$, $I_0$, the 95\% ESC at 27\% duty, and the derate/inflate factors. Chasing 2\% on S4310 $c_d$ saves tenths of a watt of aero. Putting the blade on a tabulated point with real $C_T$, or accepting that a fan is past windmill and picking another, moves tens of watts. The current $16\times12$E is already the best \emph{on-map} fan in the shortlist.

\subsection{Do not invent operating points}

If the map does not list a gentle cruise RPM at that speed, the airplane cannot have that cruise. Similarity past $J_\mathrm{windmill}$ is not a loiter mode. This is the single most expensive lesson in the repo.

\subsection{Taper first, cells second, washout last}

Taper $\lambda\approx 0.60$ from 0.28 of the outboard: fewer tip cells than a rectangle, AR 21.3, $e$ at the cap, Voc-legal strings. Older searches that dropped taper to glue more cells on the outboard made the night worse. Washout on the current wing is $\approx 0^\circ$. \textbf{Do not undo a planform that already matches the ellipse in order to print more silicon that spills at noon.}

\subsection{Motors: night $I_0$, not takeoff watts}

The 10L is the brochure motor. The 12L is the night motor. Climb still clears 1.5\,m/s by a huge margin ($\sim$15\,m/s on the $16\times12$E). Catalog kv410 is 27\% ESC duty; a rewind to $\sim$200\,rpm/V would make the 95\% ESC assumption honest without changing cruise watts in the rewind model. If a future motor has lower $I_0$ at 2000--2400\,RPM on 6S, that is worth more than another 100\,W of burst rating.

\subsection{Closing moves that actually scale}

Roughly 11\,W off night bus, or a third pack (currently forbidden), or array that produces cruise power at lower sun (camber / incidence of the cells, not more noon area). At fixed mass, raising $C_{L,\max}$ (or relaxing $1.3\to 1.2$) is the polar move that actually changes the speed you fly. T-tail downwash credits, $\pi$-tail endplate credits, and flight cooling of the cells are not those moves --- and they are conservative-physics violations if taken as free gains.

\subsection{Mindset}

\begin{itemize}
  \item Prefer an honest fail over an optimistic close.
  \item When two models disagree, take the worse number and write a FLAG, do not average them toward closure.
  \item Unused MTOW is not a license to add mass ``because we can.'' Mass is night watts.
  \item Unused climb is not a license to pick a thirstier motor ``because we have thrust.''
  \item Noon solar that spills is not endurance. Morning solar that covers cruise earlier \emph{is}.
\end{itemize}

% =====================================================================
\section{What this model does not do}
% =====================================================================

New members should not assume the following are in the physics:

\begin{itemize}
  \item 3D CFD, vortex-lattice of the full airplane, or measured glide polars (Phase 6 is the hook for those).
  \item Ground effect, banked turns, gusts, or a control-law simulation.
  \item Propeller--wing interference, swirl recovery, or blade-element theory we wrote ourselves.
  \item Measured ESC efficiency vs.\ duty; measured pack DCIR; measured cell temperature in flight.
  \item Structural sizing (spars are inside the areal density, not a beam model). No aeroelasticity.
  \item Weather other than clear-sky Ineichen. No clouds, no albedo from water, no horizon masking.
  \item Takeoff roll, landing, or a 89-hour trace with degrading packs --- only a repeating 24\,h day with a refill constraint.
  \item Payload, retracts, or a third pack (locked out).
\end{itemize}

Phase 5 (\code{monte\_carlo.py}) samples irradiance, cell temperature, array efficiency, mass growth, drag, prop factors, avionics, and pack energy, and reports $P(\mathrm{closure})$ and a tornado. It does not invent new physics; it shakes the FLAGGED knobs. Phase 6 is bench/flight protocols to replace those knobs with measurements.

% =====================================================================
\section{FLAGGED assumptions (complete working list)}
% =====================================================================

Grep \texttt{FLAGGED} in \code{solar\_uav/config.py} and the component modules. The important ones, with the conservative direction:

\begin{center}
\begin{tabular}{@{}p{0.42\textwidth}p{0.50\textwidth}@{}}
\toprule
Assumption & Conservative direction \\
\midrule
Longitude $-117.606^\circ$ (El Mirage) & $\pm 0.5^\circ$ is $\sim$2 min of solar clock; not a closure driver \\
Field elevation $865\,\mathrm{m}$ + 400\,ft AGL & already thin; sea-level density would be kind \\
Stall margin 1.2 / loiter $1.3 V_s$ and 1.5\,m/s climb & operational, not physics; $1.3 V_s$ binds night watts \\
C60 bin I & worse bins make less power \\
Cell / interconnect mass 6.5 + 1.5\,g & heavier cells hurt \\
Encapsulation 0.97, wiring 0.965, IAM diffuse 0.95 & less array \\
Cold cells $10^\circ\mathrm{C}$ for Voc & fewer cells/string (already 26) \\
String $V_\mathrm{mp}$ margin 0.95 vs.\ 19.2\,V & packing, not watts \\
Pack $R_i = 45\,\mathrm{m}\Omega$, generic OCV & more $I^2R$, slightly worse SOC accounting \\
Gearbox 0.96, ESC 0.95, ESC duty 50\% & 0.95 at 27\% duty (kv410 + $16\times12$E) is \emph{optimistic} --- noted \\
Prop $T\times 0.92$, $P\times 1.08$ & unkind to maps; table-vs-similarity FLAG can still be kind \\
Wing areal $0.84\times 0.8$\,kg/m$^2$ wetted, no spar-span scaling & user build number; may still be kind on a 6\,m spar \\
Boom $0.174\,\mathrm{kg/m}$, $\varnothing 25.4\,\mathrm{mm}$ & user build number; parasite uses the OD \\
Pod wetted from CAD, $C_f$ fully turbulent & extra parasite \\
$e_0 = 0.90$ rectangle, cap 0.98 & cap prevents taper from claiming $e=1$ \\
Downwash $2C_L/(\pi\mathrm{AR})$, no T-tail credit; $z_H=0$ & extra tail download; VLM $k=1.21$ is used (worse) \\
$C_{L,\max}$ 0.90 of 2D & higher $V_s$, higher $V_{\min}$ \\
CG at $c/4$, SM target 0.12 / min 0.08 & CG-aft would help trim; we do not take it \\
Elevator $\tau=\sqrt{c_e/c}$ & control check, not polar \\
$V_H=0.50$ (actual 0.69 from packing), $V_V=0.030$, no fin-area floor, $\mathrm{AR}_V=1.5$ & geometry / mass / yaw; no $\pi$-endplate credit \\
$T_\mathrm{amb}$ 24--40\,$^\circ\mathrm{C}$, static 75\,$^\circ\mathrm{C}$ cells & hot, no flight cooling \\
Ineichen $\pm 3\%$ 1-sigma & Phase 5 samples this \\
\bottomrule
\end{tabular}
\end{center}

If you cannot tell which way an effect goes, pick the way that makes night bus higher or SOC worse, and say so in a FLAG. Do not silently undo a conservative factor to help a candidate close.

% =====================================================================
\section{Code map}
% =====================================================================

\begin{center}
\begin{tabular}{@{}ll@{}}
\toprule
File & Role \\
\midrule
\code{solar\_uav/config.py} & locked constants, grids, FLAGGED notes \\
\code{solar\_uav/environment.py} & pvlib sun, IAM, design-day $T$, NOCT cells \\
\code{solar\_uav/asb\_physics.py} & VLM / LL / AeroBuildup / mesh inertia / XFoil envelope \\
\code{solar\_uav/aircraft.py} & geometry (independent $\ell_t$, $\ell_v$), mass, polar, trim, stall, SM, roll/yaw \\
\code{solar\_uav/lifting\_line.py} & Fourier LL, used only if washout $>0$ \\
\code{solar\_uav/mission.py} & 24\,h march, margin, closure \\
\code{solar\_uav/optimize.py} & continuous DE (includes $\ell_v$), legacy grid, inner prop loop \\
\code{solar\_uav/optimize\_kv.py} & free-$K_v$ rewind $\times$ tabulated prop at fixed airframe \\
\code{solar\_uav/cad.py} / \code{cad\_server.py} & OpenSCAD + Three.js viewer on port 8099 \\
\code{solar\_uav/components/airfoils.py} & NeuralFoil polars \\
\code{solar\_uav/components/solar\_array.py} & C60, Voc cap, bay packing \\
\code{solar\_uav/components/battery.py} & GOLD V1 SOC, 6\,A charge, $I^2R$ \\
\code{solar\_uav/components/propeller.py} & APC maps, similarity, no $n^2$-down \\
\code{solar\_uav/components/mejzlik.py} & Mejzlik maps \\
\code{solar\_uav/components/motor.py} & $K_v, R_i, I_0$, rewind model \\
\code{solar\_uav/components/propulsion.py} & matching, shortlist, climb thrust \\
\code{solar\_uav/monte\_carlo.py} & Phase 5 uncertainty \\
\code{solar\_uav/validation.py} & Phase 6 bench hooks \\
\code{.cursor/rules/conservative-physics.mdc} & the physics rule agents must follow \\
\bottomrule
\end{tabular}
\end{center}

\begin{verbatim}
.venv/bin/python scripts/validate_all.py    # Phases 1-2
.venv/bin/python scripts/eval_stability.py  # aileron / fin / loiter checks
.venv/bin/python scripts/run_asb_audit.py   # AeroSandbox vs handbook envelope
.venv/bin/python scripts/run_phase34.py     # Phases 3-4 + optimizer
.venv/bin/python scripts/run_phase56.py     # Phases 5-6 + CAD
.venv/bin/python scripts/visualize.py       # CAD / HTML viewer (http://127.0.0.1:8099/)
.venv/bin/python -m solar_uav.cad_server    # viewer only, if CAD already exported
\end{verbatim}

Constants that describe the current nearest-miss for the visualizer live at the bottom of \code{config.py} as \code{REF\_*}. Changing those does not re-run the search; it only changes what the CAD and Phase 5 treat as ``the airplane.''

% =====================================================================
\appendix
\section{Full propeller shortlist}
\label{app:props}
% =====================================================================

Electric 2-blade APC + Mejzlik maps that pass diameter $\ge 16\,\mathrm{in}$ and pitch ratio $0.40$--$0.85$ as of this writing ($N=51$).

\begin{center}
\begin{longtable}{@{}llrr@{}}
\toprule
Manufacturer & Name & Diameter [in] & Pitch [in] \\
\midrule
\endfirsthead
\toprule
Manufacturer & Name & Diameter [in] & Pitch [in] \\
\midrule
\endhead
APC & 16x7E(3D) & 16.0 & 7.0 \\
APC & 16x8E & 16.0 & 8.0 \\
APC & 16x10E & 16.0 & 10.0 \\
APC & 16x12E & 16.0 & 12.0 \\
APC & 17x7E & 17.0 & 7.0 \\
APC & 17x8E & 17.0 & 8.0 \\
APC & 17x10E & 17.0 & 10.0 \\
APC & 17x12E & 17.0 & 12.0 \\
APC & 18x8E & 18.0 & 8.0 \\
APC & 18x10E & 18.0 & 10.0 \\
APC & 18x12E & 18.0 & 12.0 \\
APC & 19x8E & 19.0 & 8.0 \\
APC & 19x10E & 19.0 & 10.0 \\
APC & 19x12E & 19.0 & 12.0 \\
APC & 20x8E & 20.0 & 8.0 \\
APC & 20x10E & 20.0 & 10.0 \\
APC & 20x11E & 20.0 & 11.0 \\
APC & 20x13E & 20.0 & 13.0 \\
APC & 20x15E & 20.0 & 15.0 \\
APC & 205x14E & 20.5 & 14.0 \\
APC & 21x13E & 21.0 & 13.0 \\
APC & 22x10E & 22.0 & 10.0 \\
APC & 22x12E & 22.0 & 12.0 \\
APC & 24x12E & 24.0 & 12.0 \\
APC & 25x125E & 25.0 & 12.5 \\
APC & 26x13E & 26.0 & 13.0 \\
APC & 26x15E & 26.0 & 15.0 \\
APC & 27x13E & 27.0 & 13.0 \\
Mejzlik & MJ-18x8E & 18.0 & 8.0 \\
Mejzlik & MJ-18x8EL & 18.0 & 8.0 \\
Mejzlik & MJ-18x10E & 18.0 & 10.0 \\
Mejzlik & MJ-20x8E & 20.0 & 8.0 \\
Mejzlik & MJ-20x8EL & 20.0 & 8.0 \\
Mejzlik & MJ-20x10E & 20.0 & 10.0 \\
Mejzlik & MJ-20x10EL & 20.0 & 10.0 \\
Mejzlik & MJ-20x11E & 20.0 & 11.0 \\
Mejzlik & MJ-20x11EW & 20.0 & 11.0 \\
Mejzlik & MJ-20x11EWL & 20.0 & 11.0 \\
Mejzlik & MJ-20x13E & 20.0 & 13.0 \\
Mejzlik & MJ-20x13EL & 20.0 & 13.0 \\
Mejzlik & MJ-20p5x12EW & 20.5 & 12.0 \\
Mejzlik & MJ-20p5x12ELW & 20.5 & 12.0 \\
Mejzlik & MJ-21x10EL & 21.0 & 10.0 \\
Mejzlik & MJ-21x13E & 21.0 & 13.0 \\
Mejzlik & MJ-21x13EL & 21.0 & 13.0 \\
Mejzlik & MJ-21x13p5EL & 21.0 & 13.5 \\
Mejzlik & MJ-21x14EL & 21.0 & 14.0 \\
Mejzlik & MJ-22x12EW & 22.0 & 12.0 \\
Mejzlik & MJ-24x10EL & 24.0 & 10.0 \\
Mejzlik & MJ-26x15E & 26.0 & 15.0 \\
Mejzlik & MJ-26x15EL & 26.0 & 15.0 \\
\bottomrule
\end{longtable}
\end{center}

\section{C60 bins (datasheet, unlaminated, STC)}

\begin{center}
\begin{tabular}{@{}lrrrrrr@{}}
\toprule
Bin & $P_\mathrm{mp}$ [W] & $V_\mathrm{mp}$ & $I_\mathrm{mp}$ & $V_\mathrm{oc}$ & $I_\mathrm{sc}$ & $\eta$ \\
\midrule
G & 3.34 & 0.574 & 5.83 & 0.682 & 6.24 & 0.218 \\
H & 3.38 & 0.577 & 5.87 & 0.684 & 6.26 & 0.221 \\
I (default) & 3.40 & 0.581 & 5.90 & 0.686 & 6.27 & 0.223 \\
J & 3.42 & 0.582 & 5.93 & 0.687 & 6.28 & 0.225 \\
\bottomrule
\end{tabular}
\end{center}

\section{How to compile this note}

The illustrated Overleaf pack (figures, charts, three-view) is \code{docs/overleaf/}: upload that folder, or the zip next to it, and compile \code{main.tex} with pdfLaTeX twice.

\begin{verbatim}
cd docs/overleaf
pdflatex main.tex
pdflatex main.tex
\end{verbatim}

Regenerate plots with \code{scripts/export\_overleaf\_figures.py}. Numbers are from the live model on 17 August 2026 (\code{reference\_design()} plus \code{outputs/phase4\_candidates.csv}).

\end{document}
